Work on the module $\Lambda^{\bullet}(W \oplus R_{2})$ of Proposition \claimref{synthesis.georgi-jarlskog-majorana-moment-map}, with $\beta_{i} = v_{i}\wedge w_{i} = n_{i} - \tfrac12$, so that $\mu = \tfrac18\sum_{i \le 5}\beta_{i}$, $\mu_{P} = \tfrac18(\beta_{1} + \beta_{2} + \beta_{3}) = \tfrac18\omega$ and $\mu_{R} = \tfrac18(\beta_{4} + \beta_{5})$ (clause (i) of that proposition). In terms of the Majorana generators, $a_{i}b_{i} = (v_{i} + w_{i})(v_{i} - w_{i}) = w_{i}v_{i} - v_{i}w_{i} = -2\beta_{i}$, a monomial of degree two, and for $i \ne l$ the product $\beta_{i}\beta_{l}$ is a monomial of degree four.

(i) $T^{\dagger} = w_{5}w_{4}w_{3}w_{2}w_{1}$, and since $v_{i}w_{i}$ commutes with $v_{l}$ and $w_{l}$ for $l \ne i$, $TT^{\dagger} = \prod_{i}v_{i}w_{i} = \prod_{i}n_{i}$, the projector onto the full state (the Majorana element raises the empty state $\nu^{c}$ to the full state, Proposition \claimref{standard-model.majorana-element}(1)). With $n_{i} = \tfrac12 + \beta_{i}$ and the $\beta_{i}$ commuting, $\prod_{i}(\tfrac12 + \beta_{i}) = \sum_{S}2^{-(5 - |S|)}\prod_{i \in S}\beta_{i}$, and $\prod_{i \in S}\beta_{i}$ is a single monomial of degree $2|S|$. The degree-zero part is $2^{-5} = \tfrac{1}{32}$, the degree-two part $\tfrac{1}{16}\sum_{i}\beta_{i} = \tfrac12\mu$, and the degree-four part $G_{4} = \tfrac18 e_{2}(\beta)$ with $e_{2}$ the second elementary symmetric polynomial. Since $\beta_{i}^{2} = \tfrac14$, $e_{2}(\beta) = \tfrac12\bigl((\sum\beta_{i})^{2} - \tfrac54\bigr) = \tfrac12(64\mu^{2} - \tfrac54)$, so $G_{4} = 4\mu^{2} - \tfrac{5}{64}$.

(ii) The monomials of $G_{4}$ are $\beta_{i}\beta_{l}$, $i < l$; those with one index among the pairs and one among the rungs sum to $X_{4} = \tfrac18(\beta_{1} + \beta_{2} + \beta_{3})(\beta_{4} + \beta_{5}) = \tfrac18(8\mu_{P})(8\mu_{R}) = 8\mu_{R}\mu_{P}$ (the factors commute). $\beta_{4} + \beta_{5} = N_{R} - 1 = 2T_{3R}$ commutes with the $\su(2)_{L}$ generators $v_{4}w_{5}$, $v_{5}w_{4}$, $n_{4} - n_{5}$ (they preserve $N_{R}$) and fails to commute with $v_{4}v_{5}$, which raises $N_{R}$ by two; $\omega$ commutes with everything diagonal. The pair-rung monomials of degree four span $\Lambda^{2}$ of the pair space tensored with $\Lambda^{2}$ of the rung space, which under $\su(4) \oplus \su(2)_{L} \oplus \su(2)_{R}$ is $(15,3,1) \oplus (15,1,3)$ (the bivectors of the six-dimensional pair space form the adjoint $15$, those of the four-dimensional rung space $(3,1) \oplus (1,3)$; standard), and an $\su(2)_{L}$ singlet lies in $(15,1,3)$.

(iii) By Proposition \claimref{synthesis.georgi-jarlskog-majorana-moment-map}(iv), $w_{4}\mu_{R}\mu_{P} = \tfrac{3}{256}(T_{\mathrm{GJ}} - w_{4}(B - L))$, so $w_{4}X_{4} = 8 \cdot \tfrac{3}{256}(\cdots) = \tfrac{3}{32}(T_{\mathrm{GJ}} - w_{4}(B - L))$; the mass forms $B(\psi, w_{4}X_{4}\varphi)$, with the pairing $B$ of Lemma \claimref{standard-model.invariant-pairing-five-modes}, follow from that proposition's (iv), scaled by $8$, with $T_{\mathrm{GJ}}$ and the ratio $3$ from Proposition \claimref{families.doublet-direction-module-georgi-jarlskog-clebsch}.

(iv) $G_{4}$ is a polynomial in $\mu$, so by Proposition \claimref{synthesis.georgi-jarlskog-majorana-moment-map}(ii) the symmetric part of the mass form of $w_{4}G_{4}$ is a multiple of the $10$ form; the multiple is zero. Indeed $\mu = \tfrac18(N - \tfrac52)$ with $N$ the total mode number, so $G_{4} = 4\mu^{2} - \tfrac{5}{64}$ is diagonal in the occupation basis and equals $\tfrac{5}{16}$, $-\tfrac{1}{16}$, $\tfrac{1}{16}$ on the states of the sixteen with $N = 0, 2, 4$. The pairing of Lemma \claimref{standard-model.invariant-pairing-five-modes} pairs a state with its complement, so $B(\psi, w_{4}\varphi) \ne 0$ on the sixteen needs $N(\psi) + N(\varphi) - 1 = 5$: one state has $N = 2$ and the other $N = 4$. The $10$ form $B(\psi, w_{4}\varphi)$ is symmetric on the sixteen, so the symmetric part of $B(\psi, w_{4}G_{4}\varphi)$ on such a pair is $\tfrac12\bigl(G_{4}(\psi) + G_{4}(\varphi)\bigr)B(\psi, w_{4}\varphi) = \tfrac12(-\tfrac{1}{16} + \tfrac{1}{16})B(\psi, w_{4}\varphi) = 0$.

(v) $(cT)(cT)^{\dagger} = |c|^{2}TT^{\dagger}$, and the expansion is linear, so $X_{4} \mapsto |c|^{2}X_{4}$; the moment map is quadratic, $\mu(cT) = |c|^{2}\mu(T)$, so the product of its two components scales as $|c|^{4}$. The Majorana entry $B(\nu^{c}, cT\nu^{c})$ is linear in $c$. At $|c| = 1$ the operators $X_{4}(T)$ and $8\mu_{R}(T)\mu_{P}(T)$ coincide; what (v) shows is that this operator occurs at second order in a vev along $T$, not only at fourth; which of the two is the physical composite is a reading.

The check builds the module as $32 \times 32$ rational matrices, expands $TT^{\dagger}$ in all $1024$ Majorana monomials by the trace form (each monomial a signed permutation, cross-checked against matrix products and its inverse), reassembles it, and confirms (i) to (v), including the mass forms of (iii) and (iv) against the pairing, and the scaling with $c = 3$.
